\section{Discussion}
\label{sec:discussion}

\subsection{Key Findings}

\textbf{Contracts and differential oracles measure fundamentally different properties.} The 40\% CU mass demonstrates that ``more tests'' cannot close the oracle-strength gap. Programs producing the correct output on contract-violating inputs are nonetheless semantically invalid --- they violate relational invariants, membership constraints, or structural postconditions that equality checks cannot express. This is a structural limitation of differential oracles, not a calibration issue.

\textbf{Adaptive PBT provides a model-consistent evaluation enhancement.} The 0.099--0.103 contribution range across 5 diverse models argues for adaptive PBT as a universally applicable evaluation tool requiring no model-specific tuning. The uniformity of the $\sim$10pp contribution suggests it reflects a fixed property of the contract oracle and input space geometry, not an artifact of any particular model's failure mode.

\textbf{The oracle gap is a threshold effect of contract presence.} $\rho = 0.136$ and the non-monotonic richness gradient indicate that any formal pre/post-condition assertion --- even a simple input validity check --- creates a qualitative oracle advantage over differential testing on CVT inputs. More complex contracts do not produce proportionally larger gaps; the threshold is crossed by contract presence alone.

\subsection{Honest Limitations}

\textbf{L1: CVT input scope.} The oracle-isolation gap is measured on contract-violating inputs, not general random inputs. Whether the 40\% gap persists on random input distributions is an open question; CVT inputs are the natural evaluation class for contract-violating behavior, and our scope qualifier is stated throughout.

\textbf{L2: Cross-model underpowering.} $n=5$ structurally precludes $\tau$ significance. $\Delta R^2 = 0.004$ is interpretable independently as a genuine null: model family adds negligible variance beyond capability proxies within the $n=5$ range.

\textbf{L3: ContractEval scope.} Results are scoped to 364 Python algorithmic tasks with inline assert contracts. Generalization to software engineering tasks, non-Python languages, or other contract formats is an open direction.

\textbf{L4: Richness gradient null.} $\rho = 0.136$ reflects precondition-dominated contract structure. Postcondition-only analysis may recover the gradient claim.

\subsection{Broader Impact}

This work contributes to the recognition that LLM code evaluation requires multiple complementary evaluation axes. Differential testing remains valuable for measuring functional input-output correctness; contract oracles add a specification-level layer that differential testing structurally cannot provide. Our results suggest that benchmark designers should treat oracle type as a first-class dimension alongside task difficulty, model scale, and test density.
